\section{Approach}

\subsection{Problem Formulation}

In federated learning, multiple clients collaboratively train a global large language model (LLM) $G$ without sharing their raw data. Consider a federated learning round $r$ where $K$ clients participate. Each client $i \in \{1, \ldots, K\}$ possesses a local dataset $\mathcal{D}_i$ and performs local training on the global model from the previous round, producing an updated client model $C_i^{(r)}$. These client models are then aggregated by the central server to form the new global model $G^{(r)}$ for round $r$.

Once the global model is trained, it can be deployed to generate responses. When presented with an input prompt $\mathbf{x}$ and the corresponding response $\mathbf{y} = (t_1, t_2, \ldots, t_T)$ generated by $G^{(r)}$, our objective is to determine which client's training data contributed most significantly to the generation of $\mathbf{y}$. Our goal is to produce an attribution vector $\mathcal{P} \in \mathbb{R}^K$ where $\mathcal{P}(i)$ quantifies the contribution of client $i$ to the generation of response $\mathbf{y}$. The client with the highest attribution score is identified as the primary source of the generated response.

This provenance capability enables critical applications in federated LLM systems, including trustworthiness verification, backdoor detection, and accountability tracking. Unlike classification tasks where a model predicts a single discrete output from a fixed set of classes, LLM provenance presents unique challenges. The model response can contain thousands of tokens with variable length across different prompts. Furthermore, when the global model generates a token, the output may stem from three potential sources: knowledge acquired from client $i$'s federated training, the model's pre-existing knowledge from pre-training, or a combination thereof.

\subsection{Background}

\textbf{Federated Model Formation.} Different model aggregation strategies are available in federated learning, such as FedAvg and FedProx. In our framework, we focus on FedAvg, which computes a weighted average of client model parameters. At each round $r$, the global model is formed through the aggregation:
\begin{equation}
\label{eq:fedavg-aggregation}
G^{(r)} = \sum_{i=1}^K \rho_i^{(r)} \cdot C_i^{(r)}
\end{equation}
where $\rho_i^{(r)} = \frac{|\mathcal{D}_i|}{\sum_{j=1}^K |\mathcal{D}_j|}$ represents the aggregation coefficient proportional to client $i$'s dataset size, and $\sum_{i=1}^K \rho_i^{(r)} = 1$.

\textbf{LLM Token Generation.} Once a global LLM is trained, it can process an input prompt $X$ to generate a response $Y$ through autoregressive token-by-token generation. The prompt is first tokenized into a sequence $(x_1, \ldots, x_t)$, which is then mapped to continuous representations by a token embedding layer. These embeddings are fed into the core of the model, which consists of a stack of $L$ identical Transformer blocks applied sequentially.

Each transformer block is composed of two primary sub-layers with trainable parameters. The first is a multi-head self-attention mechanism that enables the model to dynamically weigh the importance of different tokens relative to each other. This mechanism allows the model to aggregate information from relevant positions in the sequence, capturing complex contextual dependencies. The self-attention operates through multiple ``heads'' that learn different types of relationships simultaneously, and a final output projection layer combines the information gathered by all heads. The second sub-layer is a feed-forward network (MLP) that processes each token's representation individually. While self-attention handles contextual aggregation by routing information between positions, the MLP provides crucial non-linear transformations, allowing the model to store and apply learned knowledge to refine the meaning of each token before it is passed to the next block.

After the full sequence is processed through all $L$ transformer blocks, the model produces logits (scores) for the next token over the entire vocabulary $V$. For autoregressive generation using greedy decoding, the next token $x_{t+1}$ is selected by choosing the token with the highest score:
\begin{equation}
\label{eq:token-generation}
x_{t+1} = \operatorname*{arg\,max}_{v \in V} \left( G(x_1, \ldots, x_t) \right)_v
\end{equation}
This process continues iteratively. The generated token $x_{t+1}$ is appended to the input sequence, forming a new sequence $(x_1, \ldots, x_t, x_{t+1})$, which is then fed back into $G$ to generate the subsequent token $x_{t+2}$. This loop repeats until the model generates a special end-of-sequence token or reaches a maximum length.

\textbf{Key Enabling Property: Linear Decomposition.} The Federated Averaging (FedAvg) algorithm possesses a crucial mathematical property that enables provenance tracking in federated LLMs. Understanding this property is essential to our approach. At each round $r$, the weighted averaging of client model parameters creates a linear composition. Consider a single neuron or weight matrix $\theta$ within the model (suppressing round superscript for clarity). The global model's parameter can be expressed as:
\begin{equation}
\label{eq:linear-parameter}
\theta_{\text{global}} = \sum_{i=1}^N c_i \cdot \theta_i
\end{equation}
where $c_i$ represents the cumulative contribution coefficient of client $i$ after multiple rounds of aggregation (with $c_i = 0$ for clients not yet participated), and $\sum_{i=1}^N c_i = 1$.

This linear composition has profound implications for inference. When the global model processes an input $\mathbf{x}$, the computation at any layer can be decomposed as:
\begin{equation}
\label{eq:linear-decomposition}
\begin{split}
\mathbf{o}_{\text{global}} &= \theta_{\text{global}} \cdot \mathbf{x} \\
&= \left(\sum_{i=1}^N c_i \cdot \theta_i\right) \cdot \mathbf{x} \\
&= \sum_{i=1}^N c_i \cdot (\theta_i \cdot \mathbf{x}) \\
&= \sum_{i=1}^N c_i \cdot \mathbf{o}_i
\end{split}
\end{equation}
where $\mathbf{o}_i = \theta_i \cdot \mathbf{x}$ represents what client $i$'s model would compute independently on the same input. This decomposition reveals that the global model's output is literally a weighted sum of individual client outputs. The global model's computation can be viewed as an ensemble of client-specific computations, linearly combined through the aggregation weights. This mathematical structure is what makes provenance tractable: by analyzing how much each client's hypothetical computation $\mathbf{o}_i$ contributes to the final output, we can attribute predictions back to their source clients.

\subsection{Provenance Method}

Developing an effective provenance method for federated LLMs requires addressing three fundamental challenges that arise from the unique characteristics of autoregressive language generation and transformer architectures. For each challenge, we present a key insight that enables our solution, followed by the technical approach.

\textbf{Challenge 1: Variable-Length Autoregressive Generation.} Unlike classification tasks that produce a single prediction from a fixed set of classes, LLMs generate variable-length sequences through autoregressive token-by-token generation. Each token in the response $\mathbf{y} = (t_1, t_2, \ldots, t_T)$ may draw upon different aspects of the model's knowledge, potentially originating from different client training data. For example, in a 100-token response, tokens 1-30 might stem primarily from client A's knowledge, tokens 31-70 from the model's pre-training knowledge, and tokens 71-100 from client B's training data. This raises a fundamental question: how do we attribute an entire response when different portions may come from different sources?

\textbf{Insight and Solution.} The linear decomposition from FedAvg extends to each token generation step independently. For token $t_j$ at position $j$, the global model's computation decomposes as $\mathbf{o}_G^{(j)} = \sum_i \rho_i^{(r)} \cdot \mathbf{o}_{C_i}^{(j)}$, where each client's contribution can be measured independently. We therefore compute per-token provenance scores $\mathcal{P}_{t_j}(i)$ for each token $t_j \in \mathbf{y}$, which we then aggregate to obtain sequence-level attribution. The key is that we can measure client contributions separately at each generation step and then combine them. We aggregate these per-token contributions through summation to produce the final sequence-level provenance:
\begin{equation}
\label{eq:sequence-provenance}
\mathcal{P}(i) = \sum_{j=1}^T \mathcal{P}_{t_j}(i)
\end{equation}
This approach respects the autoregressive nature of LLM generation while providing interpretable attribution for the entire response. Tokens for which client $i$ contributed significantly will have larger $\mathcal{P}_{t_j}(i)$ values, and these accumulate to indicate overall contribution to the response.

\textbf{Challenge 2: Computational Tractability and Knowledge Localization.} Given the linear decomposition from FedAvg, we could in principle measure client contributions at every neuron in every layer of the model. However, this approach faces two critical problems. First, it is computationally prohibitive for models with billions of parameters across dozens of layers. Tracking provenance for every parameter across $K$ clients and repeating this computation for each generated token is intractable. Second, and more fundamentally, measuring all neurons would conflate relevant and irrelevant contributions, introducing significant noise into the attribution. Not all layers in a transformer LLM contribute equally to client-specific knowledge, and measuring everywhere would dilute the signal.

\textbf{Insight and Solution.} Modern transformer architectures organize computation hierarchically across layers, with different layers encoding different types of information. Research on neural network interpretability has revealed that early layers (bottom third) capture low-level linguistic features such as syntax, grammar, and basic semantic patterns. Middle layers encode compositional semantics and intermediate representations. The final layers (top third) contain more fine-grained, high-level concepts and task-specific knowledge. This hierarchical organization has direct implications for provenance: measuring provenance in late layers provides the strongest signal for attribution, as these layers encode the client-specific knowledge that most directly influences the final output.

Within each transformer block, the architecture further separates functionality between attention mechanisms and feed-forward MLP sub-layers. Recent work has demonstrated that MLP blocks function as key-value memory stores that encode factual knowledge and associations learned during training. The self-attention mechanisms, in contrast, primarily perform information routing and contextual aggregation, with less direct storage of factual knowledge. This distinction suggests that MLP outputs in late layers should be prioritized for provenance measurement, as they directly represent the stored knowledge that differentiates clients.

Based on these insights, we focus provenance computation on the final two to three transformer blocks, specifically targeting MLP projection layers (gate, up, and down projections), the attention output projection layer, and the language modeling head. For a model with $L$ layers, we select layers from the set $\mathcal{L} = \{\ell : \ell \geq L - 3\}$. This selective approach reduces computational overhead by over 80\% compared to measuring all layers, while concentrating on the components that carry the strongest client-specific signals. The attention output projection is included because, while attention performs routing, the output projection integrates multi-head results and can exhibit client-specific patterns when clients have distinct context-handling styles. The language modeling head is critical as it directly determines token predictions and carries strong provenance signals as the final transformation before vocabulary logits.

\textbf{Challenge 3: Relevance Filtering.} Even when focusing on late-layer neurons in the components we identified, a fundamental challenge remains: not all neurons within these layers are relevant for generating a particular token. A client might have large activations in neurons that are completely irrelevant to the current token prediction. For instance, neurons encoding medical knowledge would have high activations regardless of whether we are generating a medical term or a common word like ``the''. The naive approach of directly computing $\mathbf{o}_{C_i} = \theta_{C_i}^{(r)} \cdot \mathbf{x}$ for each client would conflate relevant and irrelevant activations, producing noisy attribution. We need a mechanism to automatically identify which neurons matter for the specific token being generated and weight client contributions accordingly.

\textbf{Insight and Solution.} Our solution leverages the gradient of the token prediction with respect to layer activations as an importance weighting mechanism. For a specific layer $\ell$ and token $t_j$, we compute the gradient:
\begin{equation}
\label{eq:gradient}
\mathbf{g}^\ell = \frac{\partial \log p(t_j | \mathbf{x}_{<j}; G^{(r)})}{\partial \mathbf{a}_G^\ell}
\end{equation}
where $\mathbf{a}_G^\ell$ is the activation of layer $\ell$ in the global model. This gradient quantifies how much each dimension of the layer's activation influences the final token prediction. Dimensions with larger gradient magnitudes are more influential in determining the output, while dimensions with zero or near-zero gradients are irrelevant regardless of their activation magnitude.

The key insight is that the inner product $\langle \mathbf{a}_{C_i}^\ell, \mathbf{g}^\ell \rangle$ automatically performs relevance-weighted attribution. For each dimension $k$ in the activation space, the contribution is $a_{C_i,k}^\ell \cdot g_k^\ell$. When gradient $g_k^\ell$ is zero, the contribution is zero regardless of how large $a_{C_i,k}^\ell$ might be. Conversely, when $g_k^\ell$ is large, the client's activation at that dimension directly influences the provenance score. This selective weighting focuses attribution on task-relevant computations, filtering out irrelevant neurons that would otherwise inject noise into the attribution. The mathematical intuition is that $\mathbf{g}^\ell$ defines an ``importance direction'' in the activation space---the direction most responsible for producing token $t_j$. The client whose activation $\mathbf{a}_{C_i}^\ell$ aligns most with this direction is the one whose knowledge most influenced that token generation.

\subsection{Complete Formulation}

Having presented the key challenges and insights, we now specify the complete provenance method with precise mathematical formulation. Our method operates during the autoregressive token generation process, computing provenance scores by analyzing the relationship between client-specific model activations and the gradient of the output with respect to those activations.

For the generation of a single token $t_j$ at position $j$, let $\mathbf{x}_{<j} = (t_1, \ldots, t_{j-1})$ denote the context (previous tokens). The global model computes the next token distribution as $p(t_j | \mathbf{x}_{<j}; \theta_{\text{global}})$ and selects $t_j = \arg\max p(t_j | \mathbf{x}_{<j}; \theta_{\text{global}})$. For a specific layer $\ell$ in the model, let $\mathbf{a}_{\text{global}}^\ell \in \mathbb{R}^d$ denote the activation (output) of layer $\ell$ when processing input $\mathbf{x}_{<j}$ through the global model, and let $\mathbf{g}^\ell = \frac{\partial \log p(t_j | \mathbf{x}_{<j}; \theta_{\text{global}})}{\partial \mathbf{a}_{\text{global}}^\ell} \in \mathbb{R}^d$ denote the gradient as defined in Equation~\ref{eq:gradient}.

For each client $i$, we compute what that layer would output if the client's model weights were used instead of the global weights, while keeping the input activation from the global model's forward pass. Specifically, let $\mathbf{a}_i^\ell$ denote the output of layer $\ell$ using client $i$'s parameters $\theta_i^\ell$ on the global model's input to that layer. This represents the client-specific activation pattern for the same context. The provenance score of client $i$ at layer $\ell$ for token $t_j$ is computed as the inner product between the client's activation and the output gradient:
\begin{equation}
\label{eq:layer-provenance}
S_i^{\ell}(j) = \langle \mathbf{a}_i^\ell, \mathbf{g}^\ell \rangle = \sum_{k=1}^d a_{i,k}^\ell \cdot g_k^\ell
\end{equation}

To obtain a comprehensive provenance score, we aggregate contributions across the selected layers $\mathcal{L} = \{\ell : \ell \geq L-3\}$. The total provenance score for client $i$ on token $t_j$ is:
\begin{equation}
\label{eq:token-provenance}
S_i(j) = \sum_{\ell \in \mathcal{L}} S_i^\ell(j) = \sum_{\ell \in \mathcal{L}} \langle \mathbf{a}_i^\ell, \mathbf{g}^\ell \rangle
\end{equation}
This summation accumulates evidence from different layers, with each layer capturing different aspects of the model's computation.

For a complete generated sequence $\mathbf{y} = (t_1, \ldots, t_T)$, we aggregate the per-token provenance scores to obtain an overall attribution for the entire response:
\begin{equation}
\label{eq:sequence-aggregation}
S_i = \sum_{j=1}^T S_i(j) = \sum_{j=1}^T \sum_{\ell \in \mathcal{L}} \langle \mathbf{a}_i^\ell(j), \mathbf{g}^\ell(j) \rangle
\end{equation}
where $\mathbf{a}_i^\ell(j)$ and $\mathbf{g}^\ell(j)$ denote the activations and gradients at token position $j$.

Finally, we normalize these scores using softmax to obtain a probability distribution over clients:
\begin{equation}
\label{eq:provenance-probability}
P_i = \frac{\exp(S_i)}{\sum_{k=1}^K \exp(S_k)}
\end{equation}
The client with the highest probability is identified as the primary source of the generated response:
\begin{equation}
\label{eq:attribution}
\hat{i} = \operatorname*{arg\,max}_{i} P_i
\end{equation}
This attribution provides explainability for federated model outputs and enables detection of malicious contributors in adversarial scenarios. The provenance scores $P_i$ can also be used to understand the degree of contribution from each client, providing fine-grained insights into the federated model's behavior.
